% ============================================================
% Inhalt: Mehr als zwei Zustände --- wie ein Qubit codiert
% Superposition, Register, Verschränkung und Qudits
% im FFGFT-/T0-Rahmen
% Neufassung März 2026
% Aufbaut auf Kapitel 173 (Resonator-Grundlagen)
%           und Kapitel 174 (Spin-Manipulation, Spin-Auslese)
% Shor-Algorithmus: siehe Kapitel 176
% ============================================================

% Dirac-Notation (\ket nicht im neuen Preamble definiert)
\providecommand{\ket}[1]{\left|#1\right\rangle}
\providecommand{\bra}[1]{\left\langle#1\right|}
\providecommand{\braket}[2]{\left\langle#1\middle|#2\right\rangle}

% ============================================================
\section{Die eigentliche Frage: Was steckt zwischen \texorpdfstring{$\ket{0}$}{|0>}
  und \texorpdfstring{$\ket{1}$}{|1>}?}
\label{sec:zwischen}
% ============================================================

Spin-up und Spin-down sind zwei Punkte.
Zwei Punkte sind eindeutig beschriftet --- ein Bit Information.

Aber ein \textbf{Qubit} ist kein Bit.
Die eigentliche Kraft des Qubits liegt nicht darin,
dass es $\ket{0}$ \emph{oder} $\ket{1}$ sein kann ---
das kann eine klassische Münze auch.
Sie liegt darin, dass es $\ket{0}$ \emph{und} $\ket{1}$
\emph{gleichzeitig} in einer Überlagerung sein kann ---
und zwar in einem kontinuierlichen Verhältnis:
\begin{equation}
  \ket{\psi} \;=\; \cos\!\tfrac{\theta}{2}\,\ket{0}
  \;+\; e^{i\phi}\sin\!\tfrac{\theta}{2}\,\ket{1}.
\end{equation}
Zwei Parameter: der Mischungswinkel $\theta \in [0,\pi]$
und die relative Phase $\phi \in [0, 2\pi)$.
Zusammen parametrieren sie eine \textbf{Kugeloberfläche} ---
die \textbf{Bloch-Kugel} (geometrische Darstellung
aller möglichen Einzel-Qubit-Zustände als Punkte
auf einer Kugel mit Einheitsradius).

Auf dieser Kugel gibt es \emph{unendlich viele Punkte} ---
nicht nur die zwei Pole $\ket{0}$ (Nordpol) und
$\ket{1}$ (Südpol), sondern den gesamten Äquator,
alle Breitenkreise, jeden Punkt dazwischen.
Jeder Punkt ist ein gültiger, physikalisch
unterscheidbarer Quantenzustand.

\begin{keypoint}[Ein Qubit kodiert nicht 2 Zustände --- es kodiert \emph{unendlich viele}]
Die zwei klassischen Pole $\ket{0}$ und $\ket{1}$
sind Sonderfälle einer kontinuierlichen Familie.
Was ein Qubit von einem klassischen Bit unterscheidet,
ist nicht die Zahl der diskreten Zustände ---
es ist die \textbf{kontinuierliche Geometrie} des Zustandsraums.

Das Auslesen wechselwirkt mit dem System --- jede Messung
koppelt an die Torsionskonfiguration und treibt sie
in den nächstgelegenen stabilen Extremzustand ($\ket{0}$ oder $\ket{1}$).
Das Ergebnis ist deterministisch --- es hängt vom genauen
Ausgangszustand ab, den wir vor der Messung nicht vollständig kennen.
Die scheinbare ``Wahrscheinlichkeit'' der Standard-QM beschreibt
diese epistemische Unkenntnis, keine fundamentale Zufälligkeit.

\textbf{Die Zustände selbst sind diskret.} Die Kontinuität
liegt im Phasenraum \emph{vor} der Messung ---
das Auslesen liefert immer einen der zwei stabilen Pole.
\end{keypoint}

% ============================================================
\section{Der zylindrische Phasenraum --- T0-Geometrie des Qubits}
\label{sec:zylinder}
% ============================================================

Die Bloch-Kugel ist die Standarddarstellung.
In der T0-Theorie (Dok.~034) gibt es eine
physikalisch direktere Alternative:
den \textbf{zylindrischen Phasenraum}.

Ein Qubit-Zustand wird beschrieben durch
drei geometrische Koordinaten:
\begin{equation}
  (z,\; r,\; \theta),
\end{equation}
wobei:
\begin{itemize}
  \item $z \in [-1, +1]$: die ``Höhe'' entlang der
    Zylinder\-achse --- beschreibt den Anteil
    der Torsionskonfiguration in Richtung $\ket{0}$ vs.\ $\ket{1}$.
    ($z = +1$: reine Torsion Typ~$\ket{0}$; $z = -1$: reine Torsion Typ~$\ket{1}$)
  \item $r \in [0, 1]$: der Radius im Querschnitt ---
    gibt die Stärke der intermediären Torsion an,
    d.\,h. wie weit der Zustand von beiden Polen entfernt ist.
    ($r = 0$: stabiler Pol; $r = 1$: maximale Zwischentorsion)
  \item $\theta \in [0, 2\pi)$: der Azimutalwinkel ---
    die \textbf{relative Phase} zwischen den zwei Basiskomponenten.
\end{itemize}

Die Beziehung zur Bloch-Kugel ist direkt:
$z = \cos\theta_\text{Bloch}$ und $r = |\sin\theta_\text{Bloch}|$.
Aber der Zylinder ist kein geometrisches Schönheitsmittel ---
er macht Gatter-Operationen zu \emph{geometrischen
Transformationen} im Realraum:

\begin{center}
\resizebox{\textwidth}{!}{%
\begin{tabular}{p{2.5cm}p{5cm}p{5cm}}
\toprule
\textbf{Gatter} & \textbf{Hilbertraum (Matrix)} &
\textbf{T0-Zylinder (Transformation)} \\
\midrule
\textbf{Hadamard H} &
$\frac{1}{\sqrt{2}}\begin{pmatrix}1&1\\1&-1\end{pmatrix}$ &
Vertausche $z \leftrightarrow r$, drehe Phase um $\pi/2$ ---
bewegt Pol auf Äquator \\[8pt]
\textbf{Phase Z} &
$\begin{pmatrix}1&0\\0&-1\end{pmatrix}$ &
Addiere $\pi$ zur Phase $\theta$ ---
reine Rotation um Zylinderachse \\[8pt]
\textbf{Bit-Flip X} &
$\begin{pmatrix}0&1\\1&0\end{pmatrix}$ &
2D-Rotation von $(z, r)$ um Winkel
$\alpha = \pi \cdot \Kfrak$ ---
fraktale Dämpfung eingebaut \\
\bottomrule
\end{tabular}
}
\end{center}

\begin{foundation}[FFGFT: Gatter als Torsionsrotationen]
Im zylindrischen T0-Formalismus ist jede Gatter-Operation
eine geometrische Transformation des Torsionsvektors.
Das Hadamard-Gatter dreht nicht eine Matrix ---
es kippt die Torsionskonfiguration von der
Längsachse auf die Transversalebene.
Das Z-Gatter dreht die Phase $\theta$ ---
es rotiert die Torsionswindung um ihre eigene Achse.

Der Faktor $\Kfrak$ im X-Gatter ist dabei
kein Korrekturfaktor der nachträglich eingefügt wird ---
er folgt direkt aus der fraktalen Raumzeitgeometrie
($\Dfrak = 2{,}999867$) und bedeutet:
jede Qubit-Rotation ist leicht kleiner als $\pi$,
weil die Raumzeit selbst minimal von der flachen
3D-Geometrie abweicht.
Das ist eine testbare Vorhersage:
Alle $\pi$-Gatter sollten systematisch um
$\Delta\alpha \approx 0{,}013\,\%$ vom idealen Wert abweichen.
\end{foundation}

% ============================================================
\section{Quantenregister --- exponentielles Wachstum durch Kombination}
\label{sec:register}
% ============================================================

Ein einzelnes Qubit hat unendlich viele Zwischenzustände
im Torsionsphasenraum, liefert beim Auslesen aber immer
einen der zwei stabilen Pole --- das Auslesen selbst
ist die Wechselwirkung die den Zustand in das nächste
stabile Extrem treibt. Die eigentliche Stärke kommt
durch \textbf{Kombination}.

% ============================================================
\subsection{Produktzustände: additive Klassik}
% ============================================================

Zwei unabhängige Qubits haben zusammen
$2 \times 2 = 4$ Basiszustände:
$\ket{00}$, $\ket{01}$, $\ket{10}$, $\ket{11}$.

Der allgemeine Zustand \emph{ohne Verschränkung}
(Produktzustand) ist:
\begin{equation}
  \ket{\psi_A} \otimes \ket{\psi_B}
  \;=\; (\alpha_A\ket{0} + \beta_A\ket{1})
         \otimes
         (\alpha_B\ket{0} + \beta_B\ket{1}).
\end{equation}
Das ist nur zwei getrennte Bloch-Kugeln ---
kein Gewinn über zwei unabhängige Qubits.
Der Zustandsraum ist \emph{additiv}: $2 + 2 = 4$ Parameter.

% ============================================================
\subsection{Verschränkte Zustände: multiplikatives Wachstum}
\label{sec:verschraenkung}
% ============================================================

Sobald die zwei Qubits \emph{verschränkt} sind,
kann der gemeinsame Zustand nicht mehr als Produkt
zweier Einzelzustände geschrieben werden.
Der allgemeine Zwei-Qubit-Zustand braucht
$4$ komplexe Amplituden (minus Normierung und
globale Phase: $2 \times 4 - 2 = 6$ reelle Parameter):
\begin{equation}
  \ket{\Psi} \;=\; \alpha\ket{00} + \beta\ket{01}
              + \gamma\ket{10} + \delta\ket{11}.
\end{equation}
Das ist \textbf{mehr} als zwei unabhängige Qubits ---
verschränkte Qubits tragen gemeinsame Information,
die in keinem der Einzelqubits steckt.

Für $n$ Qubits wächst der Hilbertraum exponentiell:
\begin{equation}
  \boxed{d(n) \;=\; 2^n}
\end{equation}
Reelle Parameter im allgemeinen Zustand: $2 \cdot 2^n - 2$.

\begin{center}
\resizebox{\textwidth}{!}{%
\begin{tabular}{rrrr}
\toprule
$n$ Qubits & Basiszustände & Reelle Parameter & Klassisches Äquivalent \\
\midrule
1  & 2            & 2         & 1 Bit \\
2  & 4            & 6         & 2 Bits \\
3  & 8            & 14        & 3 Bits \\
10 & 1.024        & 2.046     & 10 Bits \\
50 & $10^{15}$    & $2\times10^{15}$ & 50 Bits \\
300 & $10^{90}$   & $2\times10^{90}$ & $\approx$ Zahl der Atome im Universum \\
\bottomrule
\end{tabular}
}
\end{center}

\textbf{Das ist der eigentliche Quantenvorteil:}
Nicht ein einzelnes Qubit ist mächtig ---
sondern das kollektive Wachstum eines
\emph{verschränkten} Registers.

% ============================================================
\subsection{Bell-Zustände --- maximale Verschränkung in zwei Qubits}
\label{sec:bell}
% ============================================================

Die vier \textbf{Bell-Zustände} sind die
maximal verschränkten Zwei-Qubit-Zustände ---
jeweils eine vollständige orthonormale Basis
des verschränkten Unterraums:
\begin{align}
  \ket{\Phi^+} &= \tfrac{1}{\sqrt{2}}(\ket{00} + \ket{11}), \\
  \ket{\Phi^-} &= \tfrac{1}{\sqrt{2}}(\ket{00} - \ket{11}), \\
  \ket{\Psi^+} &= \tfrac{1}{\sqrt{2}}(\ket{01} + \ket{10}), \\
  \ket{\Psi^-} &= \tfrac{1}{\sqrt{2}}(\ket{01} - \ket{10}).
\end{align}
Jeder dieser Zustände hat die Eigenschaft:
Misst man Qubit~A, ist das Ergebnis von Qubit~B
sofort festgelegt --- egal wie weit entfernt.

In der T0-Theorie (Dok.~023, 035) gilt:
\begin{equation}
  E^{\text{T0}}(a, b) \;=\;
  -\cos(a - b) \cdot
  \left(1 - \xi\,\frac{(\Delta\theta)^2/\pi^2}{D_f}\right),
\end{equation}
die $\xi$-Dämpfung reduziert die CHSH-Verletzung von
$2\sqrt{2} \approx 2{,}8284$ auf
$\approx 2{,}8275$ --- ein Unterschied von
$\sim 10^{-4}$, der die Nicht-Lokalität lokal
erklärt, ohne sie zu beseitigen.

\begin{foundation}[FFGFT: Bell-Zustände als gemeinsame Torsionsknoten]
Im FFGFT-Bild ist ein Bell-Zustand keine
Fernwirkung zwischen zwei Qubits ---
er ist eine \textbf{topologisch verbundene
Torsionskonfiguration} über zwei Resonatoren.
Die zwei Resonatoren teilen einen gemeinsamen
Torsionsknoten, dessen Auflösung
bei jeder lokalen Messung beide Enden
gleichzeitig bestimmt.

Das ist kein Signal das sich mit Überlichtgeschwindigkeit
ausbreitet --- es ist ein gemeinsamer geometrischer
Zustand der bei der Präparation etabliert wurde.
Die lokale Auslesewechselwirkung an einem Ende
treibt die dortige Torsion deterministisch in einen
stabilen Pol --- und bestimmt damit zwingend auch
den Zustand des anderen Endes, weil beide denselben
Torsionsknoten teilen.
Die ``globale Konsequenz'' ist keine Fernwirkung,
sondern die geometrische Kohärenz des gemeinsamen Knotens.

Die $\xi$-Dämpfung (T0-Korrektur zu Bell-Korrelationen)
ist in diesem Bild die endliche Ausdehnung des
Torsionsknotens: der Knoten ist nicht punktförmig,
sondern hat eine charakteristische Länge
$\sim L_\xi = \xi \cdot \ell_P \cdot (1/\xi)^N$
auf der $\xi$-Skalenstufe $N \approx 8$.
\end{foundation}

% ============================================================
\section{Drei Wege zu mehr Zuständen --- ein Vergleich}
\label{sec:drei-wege}
% ============================================================

Aus den Kapiteln~173--175 ergeben sich drei
grundlegend verschiedene Strategien um mehr
als zwei Zustände zu kodieren.
Sie unterscheiden sich nicht nur technisch ---
sie unterscheiden sich in dem was sie
physikalisch bedeuten.

\begin{center}
\resizebox{\textwidth}{!}{%
\begin{tabular}{p{2.8cm}p{3cm}p{3.5cm}p{3.5cm}}
\toprule
\textbf{Strategie} & \textbf{Mehr Zustände durch} &
\textbf{Wachstum} & \textbf{FFGFT-Bedeutung} \\
\midrule
\textbf{Superposition}\newline (1 Qubit) &
Kontinuierliche Überlagerung von $\ket{0}$ und $\ket{1}$ &
$\infty$, aber 1 Bit bei Messung &
Intermediäre Torsion auf der Bloch-Kugel \\[8pt]
\textbf{Register}\newline ($n$ Qubits) &
Verschränkung mehrerer Qubits &
$2^n$ (exponentiell) &
$2^n$-dimensionaler Torsions-Produktraum; verschränkt = topologisch verbundene Knoten \\[8pt]
\textbf{Qudit}\newline ($d > 2$, 1 Objekt) &
Höhere Energieniveaus, Orbitalzustände, Exziton-Zahl &
$d$ (linear) pro Objekt, $d^n$ für $n$ Qudits &
$d$ stabile Torsionskonfigurationen in einem Resonator \\
\bottomrule
\end{tabular}
}
\end{center}

\begin{keypoint}[Keine Strategie dominiert in allen Bereichen]
Register-Verschränkung skaliert am stärksten ($2^n$),
braucht aber $n$ physische Objekte und fehlerkorrigierte
Verschränkungsoperationen zwischen ihnen.

Qudit-Kodierung gewinnt Information \emph{pro Objekt}
($\log_2 d > 1$), aber das exponentielle
Wachstum ist langsamer ($d^n < 2^n$ für $d < 2^n$... 
warte: für $d=4$ und 1 Objekt: 2 Bits vs. 2 Qubits: 2 Bits ---
dasselbe! Aber 2 Ququarts ($d=4$) = $4^2=16$
vs. 4 Qubits = $2^4=16$ --- auch dasselbe.
Qudits und Qubit-Register sind \emph{mathematisch äquivalent},
physikalisch aber unterschiedlich robust).

Superposition ist immer da --- sie ist die Grundlage
für Quanteninterferenz, die eigentliche Triebkraft
von Quantenalgorithmen.
\end{keypoint}

% ============================================================
\section{Quanteninterferenz --- warum Superposition nützlich ist}
\label{sec:interferenz}
% ============================================================

Superposition allein wäre eine bloße Überlagerung ---
das Auslesen würde immer einen der zwei stabilen Pole
liefern, deterministisch abhängig vom genauen Phasenzustand.

Was Superposition zur \emph{Ressource} macht ist die
\textbf{Interferenz}: Amplituden können sich
konstruktiv (verstärken) oder destruktiv (auslöschen)
überlagern --- je nach Phase $\phi$.

Das ist das Prinzip hinter allen nützlichen
Quantenalgorithmen:
\begin{enumerate}
  \item Präpariere das Register in einer
    gleichmäßigen Zwischentorsion über alle Eingaben
    (mit dem Hadamard-Gatter).
  \item Wende eine Funktion an, die die Phasen
    der richtigen Antworten erhöht und die
    der falschen erniedrigt.
  \item Lass Interferenz die falschen Antworten
    in instabile Torsionskonfigurationen treiben
    und die richtige in eine stabile.
  \item Das Auslesen treibt das System deterministisch
    in den nächsten stabilen Pol ---
    der nach der Interferenz mit überwältigender
    Häufigkeit der gesuchten Antwort entspricht.
\end{enumerate}

Ohne Interferenz: Monte-Carlo (klassisch raten).
Mit Interferenz: Grover ($\sqrt{N}$ statt $N$ Schritte),
Shor (Faktorisierung in polynomialer Zeit).

\begin{foundation}[FFGFT: Interferenz als Torsionsphasenüberlagerung]
In FFGFT ist die Phase $\phi$ keine abstrakte
komplexe Zahl --- sie ist die Rotationsstellung
der intermediären Torsionskonfiguration (Kapitel~174).

Konstruktive Interferenz: Zwei Torsionsphasen
liegen gleich --- sie verstärken sich zu einer
geometrisch stabileren Konfiguration.
Destruktive Interferenz: Zwei entgegengesetzte
Torsionsphasen löschen sich aus ---
die resultierende Konfiguration hat keinen
stabilen geometrischen Anker und wird bei der
Auslesewechselwirkung nicht angesteuert.

Ein Quantenalgorithmus ist in FFGFT-Sprache:
eine Folge von Torsionsrotationen, die die
Phasenkonfiguration des Registers so ausrichtet,
dass die gesuchte Antwort in einer stabilen
Torsionskonfiguration endet --- alle anderen
in geometrisch instabilen.
Das Auslesen ist dann keine zufällige Auswahl,
sondern eine deterministische Reaktion des
Resonators auf die Auslesewechselwirkung.
\end{foundation}

% ============================================================
\section{QND-Messung --- Auslesen ohne Zerstören}
\label{sec:qnd}
% ============================================================

Im vorigen Abschnitt wurde die Auslesewechselwirkung
als der Schritt beschrieben, der die intermediäre
Torsionskonfiguration in einen stabilen Pol treibt.
Das ist korrekt --- für \emph{invasive} Messverfahren.

Es gibt aber eine Klasse von Ausleseverfahren die
diesen Schritt nicht vollzieht:
die \textbf{Quantum-Non-Demolition-Messung} (QND).

% ============================================================
\subsection{Was QND bedeutet}
% ============================================================

Eine QND-Messung liest eine Observable $\hat{A}$ aus,
ohne den dazugehörigen Eigenzustand zu verändern.
Die formale Bedingung:
\begin{equation}
  [\hat{A},\; \hat{H}_\text{Messung}] \;=\; 0.
\end{equation}
Der Messoperator $\hat{H}_\text{Messung}$ kommutiert
mit der gemessenen Observable --- er verändert den
Eigenwert nicht, weil er denselben Eigenraum respektiert.

Praktisch bedeutet das: derselbe Zustand kann
\textbf{wiederholt gemessen} werden und liefert
jedes Mal dasselbe Ergebnis --- weil die Messung
den Zustand selbst nicht berührt.

Das wichtigste Beispiel aus Kapitel~174:
Die \textbf{Faraday/Kerr-Rotation}.
Ein linear polarisierter Laserstrahl passiert den
Quantenpunkt ohne absorbiert zu werden ---
er dreht seine Polarisationsebene um $\theta_F \propto S_z$.
Kein Photon wird absorbiert. Kein Energieübertrag
in die Spin-Mode. Der Spinzustand vor und nach
der Messung ist identisch.

% ============================================================
\subsection{QND als empirischer Beweis für Determinismus}
\label{sec:qnd-beweis}
% ============================================================

Die QND-Messung ist aus FFGFT-Sicht weit mehr
als eine technisch nützliche Eigenschaft ---
sie ist ein \textbf{empirischer Beweis}
für die Realität stabiler Torsionskonfigurationen.

Das Argument ist direkt:

\begin{enumerate}
  \item Die Faraday-Rotation misst $S_z$ ohne den
    Spinzustand zu verändern.
  \item Wiederholte Messungen liefern immer dasselbe Ergebnis.
  \item Also hatte der Zustand diesen Wert
    \emph{vor jeder Messung} --- und zwar definit,
    nicht als abstrakte Superposition die erst durch
    Messung ``Realität wird''.
  \item Der Zustand war real bevor wir hingeschaut haben.
\end{enumerate}

\begin{foundation}[FFGFT: QND-Messung und stabile Torsionskonfiguration]
In FFGFT ist die Faraday-Rotation deshalb zustandserhaltend,
weil eine stabile Torsionskonfiguration
(Spin-up = topologisch fixierte Windung in eine Richtung)
durch eine dispersive, nicht-resonante Kopplung
nicht aus ihrem Zustand getrieben werden kann.

Das Photon ``sieht'' die Torsion --- es dreht seine
Polarisation entsprechend --- aber es überträgt keine
Energie in die Torsionsmode, weil es nicht resonant ist.
Der geometrische Zustand ist invariant unter dieser Kopplung.

Das ist der Unterschied zur invasiven Messung:
\begin{center}
\begin{tabular}{p{3cm}p{5cm}p{5cm}}
\toprule
\textbf{Verfahren} & \textbf{Wechselwirkung} & \textbf{Zustand danach} \\
\midrule
Resonante Fluoreszenz &
Photon wird absorbiert, regt Übergang an &
Zustand verändert (angeregt, dann Zerfall) \\[6pt]
Spin-zu-Ladungs-\newline Konversion &
Elektron tunnelt aus dem Resonator &
Zustand zerstört (Elektron weg) \\[6pt]
\textbf{Faraday/Kerr} &
Photon passiert dispersiv, kein Absorptions &
\textbf{Zustand unverändert --- QND} \\
\bottomrule
\end{tabular}
\end{center}

\textbf{Die philosophische Konsequenz ist eindeutig:}
Wenn ein Zustand wiederholt ohne Störung gemessen
werden kann, dann ist er nicht durch die Messung
``erzeugt'' worden --- er war vorher da.
Die Standard-QM-Interpretation
(``der Zustand existiert erst durch Messung'')
ist mit QND-Messungen nicht vereinbar.

In FFGFT: Die Torsionskonfiguration ist eine
geometrische Realität --- stabil, definit, präexistent.
Die Messung liest sie ab. Sie erschafft sie nicht.
\end{foundation}

% ============================================================
\subsection{Grenzen der QND-Messung}
% ============================================================

QND ist möglich für Observablen die mit dem
Messoperator kommutieren --- typisch: $S_z$ (Spin-Komponente
entlang der Quantisierungsachse).

Was QND \emph{nicht} leisten kann:
Eine intermediäre Torsionskonfiguration
(Superposition $\alpha\ket{\uparrow} + \beta\ket{\downarrow}$,
also ein Zustand der \emph{kein} $S_z$-Eigenzustand ist)
wird durch eine $S_z$-QND-Messung in einen Eigenzustand
getrieben --- deterministisch, abhängig vom genauen
Ausgangszustand, aber die Zwischentorsion ist danach weg.

Warum das so ist, zeigt das Potentiallandschafts-Bild.

% ============================================================
\subsection{Potentiallandschaft: stabile Minima, Übergänge und Auslesestärke}
\label{sec:potential}
% ============================================================

\begin{foundation}[FFGFT: Torsionspotential, Übergänge und Auslesestärke]

Das vollständige Bild der Qubit-Zustände ist eine
\textbf{Potentiallandschaft mit zwei stabilen Minima}
und einem Sattel dazwischen:

\medskip
\begin{center}
\begin{tabular}{p{3.5cm}p{9cm}}
\toprule
\textbf{Region} & \textbf{Physikalische Bedeutung} \\
\midrule
Minimum $\ket{\uparrow}$ &
Stabile Torsionskonfiguration Typ~1.
Tief genug dass schwache Kopplungen den Zustand
nicht heraus\-heben --- QND-Auslese möglich. \\[6pt]
Minimum $\ket{\downarrow}$ &
Stabile Torsionskonfiguration Typ~2.
Dieselbe Tiefe, spiegelbildlich. \\[6pt]
Abhang / Sattel &
Intermediäre Torsion: die Phase $\phi$ bestimmt
wie weit der Zustand von welchem Minimum entfernt ist.
Kein echtes Minimum --- metastabil auf der
Kohärenzzeit $T_2$.\\[6pt]
Übergang (Phase überlagert) &
Jeder Punkt zwischen den zwei Minima ist ein
geometrisch definierter Zustand ---
kontinuierlich parametriert durch $(\theta, \phi)$
auf der Bloch-Kugel, diskret in den Endpunkten. \\
\bottomrule
\end{tabular}
\end{center}
\medskip

\textbf{Die Auslesestärke entscheidet was passiert:}

\begin{center}
\begin{tabular}{p{3cm}p{4.5cm}p{5cm}}
\toprule
\textbf{Ausgangszustand} & \textbf{Schwache Kopplung} & \textbf{Starke Kopplung} \\
\midrule
Stabiles Minimum\newline ($\ket{\uparrow}$ oder $\ket{\downarrow}$) &
Kopplung hebt Zustand nicht aus dem Minimum ---
\textbf{QND}: Zustand bleibt, Ergebnis reproduzierbar &
Kopplung kann Zustand aus dem Minimum heben ---
invasiv, Zustand verändert \\[8pt]
Abhang / Sattel\newline (Superposition) &
Auch schwache Kopplung reicht aus den Zustand
zu treiben --- er ``rollt'' in das nächste Minimum.
\textbf{Welches}: bestimmt durch die genaue Phase $\phi$
(epistemisch unbekannt, deterministisch) &
Dasselbe, nur schneller ---
der Zustand springt sofort in das nächste Minimum \\
\bottomrule
\end{tabular}
\end{center}
\medskip

\textbf{Die entscheidende Einsicht:}
Es gibt keine prinzipielle Grenze zwischen
``QND-möglich'' und ``QND-unmöglich'' ---
es gibt nur \textbf{stabilere und weniger stabile Zustände}
in einer kontinuierlichen Potentiallandschaft.

Ein Zustand im tiefen Minimum ist stabil gegenüber
der Auslesekopplung --- das Minimum muss überwunden werden.
Ein Zustand auf dem Abhang ist instabil ---
jede Kopplung, egal wie schwach, reicht aus ihn
ins nächste Minimum zu treiben.

Und der ``Sprung'' in die Endposition ist nicht zufällig ---
er ist deterministisch durch die Phase des Ausgangszustands
auf dem Abhang bestimmt:
ein Zustand nahe $\ket{\uparrow}$ landet in $\ket{\uparrow}$,
ein Zustand nahe $\ket{\downarrow}$ landet in $\ket{\downarrow}$,
ein Zustand auf dem Sattel ($\theta = \pi/2$) landet
abhängig von der feinen Phase $\phi$.

Die scheinbare ``Zufälligkeit'' der Standard-QM
ist die epistemische Unkenntnis dieser Phase ---
kein ontologischer Zufall des Systems.

\textbf{In FFGFT-Sprache:}
Stabile Torsionskonfigurationen sind geometrische Minima.
Intermediäre Torsionen sind Positionen auf dem Abhang.
Die Auslesewechselwirkung ist eine externe Kraft
die das System die Potentiallandschaft entlangtreibt.
Je nach Startposition und Stärke der Kraft:
entweder Verbleib im Minimum (QND)
oder deterministisches Rollen in das nächste (invasiv).

\textbf{Das Bild gilt für alle $d$ Niveaus eines Qudits:}
Nicht zwei Minima, sondern $d$ Minima ---
mit $d-1$ Sattelpunkten zwischen ihnen.
Die Übergänge zwischen benachbarten Niveaus
folgen derselben Logik: stabil im Minimum,
deterministisch getrieben vom Abhang.
\end{foundation}

\begin{keypoint}[QND und Superposition --- kein Widerspruch]
Ein stabiler Eigenzustand ($\ket{\uparrow}$ oder $\ket{\downarrow}$):
sitzt im tiefen Minimum --- schwache Kopplung liest aus
ohne zu treiben. QND möglich. Ergebnis reproduzierbar.

Eine intermediäre Torsion ($\alpha\ket{\uparrow}+\beta\ket{\downarrow}$):
sitzt auf dem Abhang --- jede Kopplung treibt
in das nächste Minimum. Deterministisch, nicht zufällig.
Die Phase $\phi$ bestimmt welches Minimum erreicht wird.

Die QND-Messung \emph{beweist} die Realität stabiler Pole.
Die verändernde Messung \emph{beweist} die Instabilität
intermediärer Torsionen --- und die Determinismus-Struktur
der Potentiallandschaft.
Zusammen: vollständig deterministisches,
geometrisch fundiertes Bild.
\end{keypoint}

% ============================================================
\section{Zwei Ebenen der \texorpdfstring{$\xi$}{xi}-Hierarchie:
  Spin-Qubit und photonischer Resonator}
\label{sec:zwei-ebenen}
% ============================================================

Die bisherigen Abschnitte haben Quantenzustände
hauptsächlich am Beispiel des Spin-Qubits entwickelt ---
einem einzelnen Elektron in einem Nanometer-Käfig.
Es gibt aber eine qualitativ andere Realisierungsklasse:
den \textbf{photonischen Resonator}.

Beide sind Resonatoren im Sinne von Kapitel~173.
Beide erzwingen diskrete Zustände durch Geometrie.
Aber sie sitzen auf \emph{verschiedenen Stufen}
der $\xi$-Hierarchie --- und das hat fundamentale
Konsequenzen für Temperatur, Stabilität und
die Art der Information die sie tragen.

% ============================================================
\subsection{Spin-Qubit: minimaler Resonator auf Stufe \texorpdfstring{$N \approx 8$}{N≈8}}
% ============================================================

Ein Spin-Qubit ist das absolute Minimum eines
Quantenresonators (Kapitel~173):
ein einzelnes Elektron, ein einziges Energieniveau,
zwei erlaubte Orientierungen.

Die Energieskala kommt aus dem
\textbf{Quantum-Confinement} auf der $\xi$-Stufe $N \approx 8$
(Nanometer-Skala, Abschnitt~\ref{sec:zylinder}).
Der Abstand zwischen den zwei Spin-Zuständen
in einem typischen Magnetfeld ($B \approx 1$\,T)
liegt bei:
\begin{equation}
  \Delta E_Z \;=\; g^* \mu_B B \;\approx\; 0{,}1\text{--}1\,\text{meV}.
\end{equation}
Die thermische Energie bei Raumtemperatur ist:
\begin{equation}
  k_B T(300\,\text{K}) \;=\; 26\,\text{meV}.
\end{equation}
Also $\Delta E_Z \ll k_BT$ --- thermisches Rauschen
wirft das System ständig aus dem Minimum heraus.
Tiefe Temperaturen ($T < 1$\,K, $k_BT < 0{,}1$\,meV)
sind nicht optional, sondern zwingend ---
nicht um Superposition zu erhalten,
sondern schlicht um in einem der zwei Minima
zu \emph{bleiben}.

% ============================================================
\subsection{Photonischer Resonator: kollektive Geometrie auf Stufe
  \texorpdfstring{$N \approx 9$--10}{N≈9-10}}
% ============================================================

Ein Wellenleiter aus Lithiumniobat (LiNbO$_3$),
ein Mach-Zehnder-Interferometer, ein Ringresonator ---
das sind Resonatoren aus $10^{18}$ bis $10^{23}$ Atomen
in perfekter Gitterordnung.

Der entscheidende Unterschied zum Spin-Qubit:
\textbf{Der Resonator ist nicht ein einzelnes Teilchen ---
er ist die makroskopische Geometrie des gesamten Kristalls.}

Das Photon das darin propagiert gehört nicht einem
einzelnen Atom. Es gehört der kollektiven Mode
des gesamten Resonators.
Die relevante Diskretheit kommt nicht aus
Quantum-Confinement eines einzelnen Elektrons,
sondern aus den \textbf{Randbedingungen der
makroskopischen Geometrie} --- Länge, Brechungsindex,
Spiegelreflektivität.

Die Energieskala eines optischen Photons:
\begin{equation}
  E_\text{Photon} \;\approx\; 1\text{--}2\,\text{eV}
  \quad (780\text{--}1550\,\text{nm}),
\end{equation}
also:
\begin{equation}
  \frac{E_\text{Photon}}{k_BT(300\,\text{K})}
  \;\approx\; \frac{1\,\text{eV}}{26\,\text{meV}}
  \;\approx\; 40.
\end{equation}
Thermisches Rauschen kann spontan \emph{keine} optischen
Photonen erzeugen --- das System sitzt weit unterhalb
der thermischen Schwelle für optische Anregungen.
Bedingung~2 aus Kapitel~173 ($\Delta E \gg k_BT$)
ist um einen Faktor~40 übererfüllt.

\begin{important}[Warum Photonik bei Raumtemperatur funktioniert]
Nicht weil das photonische System robuster gebaut ist ---
sondern weil es auf einer höheren $\xi$-Stufe
der Energiehierarchie operiert.

Die \textbf{geometrische Robustheit} (Bedingung~1)
ist durch den makroskopischen Einkristall mit
$\sim 10^{23}$ Gitterplätzen in perfekter Ordnung
außerordentlich gut erfüllt ---
viel robuster als ein einzelner Quantenpunkt.
Einzelne Phononen (thermische Gitterschwingungen)
stören die kollektive Mode kaum.

Die \textbf{thermische Robustheit} (Bedingung~2)
ist durch die eV-Energieskala automatisch gegeben ---
thermische Photonen im Infrarot haben $\sim 25$\,meV,
optische Photonen haben $\sim 1$\,eV:
vier Größenordnungen Sicherheitsabstand.
\end{important}

% ============================================================
\subsection{Nicht auf zwei Zustände beschränkt:
  unendlich viele analoge Übergänge}
% ============================================================

Das Spin-Qubit ist auf zwei Zustände reduziert ---
das ist seine Stärke (Einfachheit, Kontrolle)
und seine Schwäche (Kühlung zwingend, fragil).

Der photonische Resonator ist strukturell anders:
Er ist \emph{nicht} auf die Minimal-Zwei-Zustands-Struktur
des Spin-Up/Down beschränkt.

In einem MZI (Mach-Zehnder-Interferometer)
läuft das Photon durch zwei Pfade gleichzeitig ---
mit einer kontinuierlich einstellbaren Phase $\phi$
zwischen den Pfaden.
Die Phase ist eine analoge Größe die
$[0, 2\pi)$ kontinuierlich durchlaufen kann.
Das Interferenzergebnis ist $\cos^2(\phi/2)$ ---
eine kontinuierliche Übertragungsfunktion,
keine diskrete Zweistufigkeit.

\begin{keypoint}[Analog oder digital --- eine Wahl der Betriebsweise]
Ein photonischer Resonator kann \textbf{beide Modi} betreiben:

\textbf{Digital:} Phase auf $\phi = 0$ oder $\phi = \pi$ gesetzt ---
dann ist Transmission 1 oder 0.
Zwei diskrete Zustände, wie ein klassisches Bit.

\textbf{Analog:} Phase kontinuierlich zwischen 0 und $2\pi$ ---
dann ist Transmission eine kontinuierliche Zahl
zwischen 0 und 1. Unendlich viele Zwischenzustände,
alle bei Raumtemperatur stabil, weil die Energieskala
thermisch unerreichbar ist.

Das analoge Auswerten dieser Übergänge erfordert
\emph{keine} Kühlung --- aber es erfordert
geometrische Präzision des Resonators
(Bedingung~1 aus Kapitel~173).
\end{keypoint}

Das ist der fundamentale Unterschied zum Spin-Qubit:
Beim Spin-Qubit liegt die analoge Information
(die Phase $\phi$ auf dem Abhang der Potentiallandschaft)
in einem thermisch instabilen Zwischenzustand ---
$T_2$ bricht bei Raumtemperatur in Nanosekunden zusammen.

Beim photonischen Resonator liegt die analoge Information
in der Phase des Photons --- und Photonen sind
\textbf{thermisch stabil}, weil ihre Energieskala
weit über $k_BT$ liegt.
Der Zustand auf dem ``Abhang'' ist hier kein
fragiles Minimum --- er ist ein geometrisch
stabilisierter Modenparameter des Kristalls.

% ============================================================
\subsection{Vergleich der zwei Ebenen in der \texorpdfstring{$\xi$}{xi}-Hierarchie}
\label{sec:xi-vergleich}
% ============================================================

\begin{center}
\resizebox{\textwidth}{!}{%
\begin{tabular}{p{3.5cm}p{5cm}p{5cm}}
\toprule
& \textbf{Spin-Qubit} & \textbf{Photonischer Resonator} \\
\midrule
\textbf{$\xi$-Stufe} &
$N \approx 8$ (nm-Skala) &
$N \approx 9$--10 ($\mu$m--mm-Skala) \\[4pt]
\textbf{Resonator} &
1 Elektron im nm-Käfig &
$10^{18}$--$10^{23}$ Atome als kollektiver Kristall \\[4pt]
\textbf{Energieskala} &
$\Delta E \approx 0{,}1$--1\,meV &
$E_\text{Photon} \approx 1$--2\,eV \\[4pt]
\textbf{$\Delta E / k_BT$(300\,K)} &
$\approx 0{,}004$--0{,}04 (zu klein!) &
$\approx 40$ (robust) \\[4pt]
\textbf{Kühlung nötig?} &
Ja, $T < 1$\,K zwingend &
Nein, Raumtemperatur \\[4pt]
\textbf{Analoge Zustände} &
Fragil: $T_2 \sim$ ns bei 300\,K &
Stabil: Phase thermisch unerreichbar \\[4pt]
\textbf{Diskretheit aus} &
Quantum-Confinement (1 Teilchen) &
Makroskopische Modengeometrie \\[4pt]
\textbf{Bedingung~1} &
Einzelner QD: präzise, aber störbar &
Makroeinkristall: extrem robust \\[4pt]
\textbf{Minimale Struktur} &
2 Zustände (Spin-up/down) &
Nicht beschränkt: kontinuierliche Phase \\[4pt]
\textbf{Typische Anwendung} &
Digitales Qubit, Quantengatter &
Analoge Signalverarbeitung,\\
& & photonisches Rechnen, 6G \\
\bottomrule
\end{tabular}
}
\end{center}

\begin{foundation}[FFGFT: Zwei Stufen der $\xi$-Hierarchie ---
  dieselbe Physik, verschiedene Skalen]
Spin-Qubit und photonischer Resonator sind
in FFGFT keine grundlegend verschiedenen Systeme ---
sie sind Resonatoren auf verschiedenen Stufen
der $\xi$-Skalierungsleiter.

Auf Stufe $N \approx 8$: das Confinement-Energieminimum
liegt bei meV --- thermisch nicht robust bei 300\,K.
Die zwei stabilen Torsionskonfigurationen (Spin-up/down)
sind geometrisch erzwungen, aber thermisch überwältigbar.

Auf Stufe $N \approx 9$--10: das kollektive Modenminimum
liegt bei eV --- thermisch vollständig robust.
Die stabilen Konfigurationen (Photonenmoden, Interferenzzustände)
sind geometrisch erzwungen \emph{und} thermisch stabil.

Der $\xi$-Parameter erzeugt eine universelle
Energiehierarchie: jede Stufe trägt die Energie
um den Faktor $1/\xi \approx 7500$ höher als die vorherige.
Von meV (Stufe~8) zu eV (Stufe~9--10):
der Faktor beträgt $\sim 10^3$--$10^4$ ---
konsistent mit der $\xi$-Skalierung.

Das erklärt nicht nur Photonik bei Raumtemperatur ---
es erklärt warum biologische Systeme photonische
Mechanismen (Photosynthese, Vogelnavigation via Licht)
und keine Spin-Qubit-Systeme für Quanteneffekte nutzen:
die Natur arbeitet auf der thermisch robusten Stufe.
\end{foundation}

% ============================================================
\subsection{Die analoge Auflösungsgrenze --- und ihre Fortsetzung
  in der Hochfrequenztechnik}
\label{sec:aufloesung}
% ============================================================

Der photonische Resonator ist thermisch robust ---
aber nur solange man nicht zu fein auflöst.
Sobald man analoge Zwischenwerte unterscheiden will,
teilt man den Energieraum $E_\text{Photon}$ in
$N$ unterscheidbare Stufen.
Die Energiedifferenz zwischen benachbarten Stufen:
\begin{equation}
  \Delta E_\text{Stufe} \;=\; \frac{E_\text{Photon}}{N}.
\end{equation}
Bedingung~2 aus Kapitel~173 gilt unerbittlich:
\begin{equation}
  \Delta E_\text{Stufe} \;\gg\; k_BT
  \quad\Longrightarrow\quad
  N \;\ll\; \frac{E_\text{Photon}}{k_BT}.
\end{equation}
Bei Raumtemperatur und optischen Photonen ($E \approx 1$\,eV):
\begin{equation}
  N_\text{max}(300\,\text{K})
  \;\approx\; \frac{1\,\text{eV}}{26\,\text{meV}}
  \;\approx\; 40.
\end{equation}
Die maximale analoge Auflösung bei Raumtemperatur
beträgt etwa 40 unterscheidbare Stufen ---
entsprechend $\log_2(40) \approx 5{,}3$\,Bits
pro Photon (Shannon-Grenze für dieses Energieniveau).

Das photonische System entkommt dem thermischen Problem
nicht --- es verschiebt es auf eine feinere Skala.
Mehr Auflösung erfordert entweder Kühlung
(womit man dieselbe Argumentation wie beim Spin-Qubit erhält)
oder eine noch höhere Energieskala.

\begin{foundation}[Die gesamte Geschichte der Hochfrequenztechnik
  als Aufstieg in der $\xi$-Hierarchie]

Das ist keine neue Einsicht --- es ist die älteste
Ingenieursweisheit der Nachrichtentechnik,
jetzt geometrisch begründet.

Die Hochfrequenztechnik klettert seit hundert Jahren
dieselbe $\xi$-Hierarchie hoch.
Bei jeder Frequenzstufe verbessert sich das
Signal-Rausch-Verhältnis --- nicht weil die Ingenieure
bessere Schaltkreise bauten, sondern weil
$\Delta E / k_BT$ mit jeder Stufe wächst.

\begin{center}
\begin{tabular}{p{2.5cm}p{1.8cm}p{2.2cm}p{1.5cm}p{1.5cm}p{3cm}}
\toprule
\textbf{Technologie} &
\textbf{Frequenz} &
\textbf{$E = hf$} &
\textbf{$E/k_BT$} &
\textbf{$N$-Stufe} &
\textbf{Thermisches Regime} \\
\midrule
AM-Rundfunk &
1\,MHz &
$4\times10^{-9}$\,eV &
$1{,}5\times10^{-7}$ &
$\approx 5{,}5$ &
Thermisch vollständig dominiert;
Rauschen $\gg$ Signal \\[4pt]
UKW/FM &
100\,MHz &
$4\times10^{-7}$\,eV &
$1{,}5\times10^{-5}$ &
$\approx 6{,}4$ &
Immer noch stark rauschbegrenzt;
starke Sender nötig \\[4pt]
Mikrowelle / Radar &
10\,GHz &
$4\times10^{-5}$\,eV &
$1{,}5\times10^{-3}$ &
$\approx 7{,}3$ &
Rauschen beherrschbar;
rauscharme Verstärker nötig \\[4pt]
5G / mmWave &
30\,GHz &
$1{,}2\times10^{-4}$\,eV &
$5\times10^{-3}$ &
$\approx 7{,}5$ &
Übergangsbereich;
Kühlung empfindlichster Stufen
(Radioteleskope) \\[4pt]
THz &
1\,THz &
$4\times10^{-3}$\,eV &
$0{,}15$ &
$\approx 7{,}8$ &
Schwierige Zone: $E \sim k_BT$;
deshalb THz-Technik so aufwendig \\[4pt]
Nahinfrarot / LiDAR &
200\,THz &
$0{,}8$\,eV &
31 &
$\approx 8{,}8$ &
Thermisch robust;
Raumtemperatur problemlos \\[4pt]
Optik / Glasfaser &
300\,THz &
$1{,}2$\,eV &
46 &
$\approx 9{,}0$ &
Thermisch vollständig stabil;
$N_\text{max} \approx 46$ analoge Stufen \\[4pt]
Kurzwellige UV &
1.000\,THz &
$4$\,eV &
154 &
$\approx 9{,}4$ &
Hohe analoge Auflösung;
$N_\text{max} \approx 150$ Stufen \\
\bottomrule
\end{tabular}
\end{center}

\medskip
\textbf{Der einheitliche Zusammenhang:}

Jede Stufe in der $\xi$-Hierarchie trägt die Energieskala
um den Faktor $1/\xi \approx 7500$ höher als die vorherige.
Von $N = 6$ (AM, MHz) nach $N = 9$ (Optik, THz)
überbrücken drei Stufen den Faktor
$(1/\xi)^3 \approx 4 \times 10^{11}$ ---
und tatsächlich liegt die Frequenz der Glasfaser
$\sim 10^8$-fach über AM-Rundfunk.

\textbf{Die Glasfaser ist in diesem Bild keine neue Technologie ---
sie ist die konsequente Fortsetzung desselben Trends
auf die nächste $\xi$-Stufe.}

Und die Auflösungsgrenze gilt auf jeder Stufe:
wer bei optischen Wellenlängen mehr als $\sim$40 analoge
Stufen auflösen will, muss kühlen ---
genau wie der Radioingenieur 1930 seinen Empfänger
abschirmte und seinen Verstärker optimierte um
im thermischen Rauschen noch ein Signal zu finden.

Das Problem ist dasselbe. Die Skala verschiebt sich.
Die Physik ändert sich nicht.

\textbf{In FFGFT:}
Die universelle Auflösungsgrenze pro Temperaturstufe:
\begin{equation}
  \boxed{N_\text{max}(T, N_\xi)
  \;\approx\;
  \frac{E_0 \cdot (1/\xi)^{N_\xi - N_0}}{k_B T}}
\end{equation}
mit $E_0 = k_B T_0$ bei der Referenzstufe $N_0$.
Jede $\xi$-Stufe höher multipliziert die
erreichbare Auflösung um den Faktor $1/\xi \approx 7500$.
\end{foundation}

% ============================================================
\subsection{Oberhalb der Optik --- drei Regime bis zur Paarerzeugungs-Schwelle}
\label{sec:oberhalb-optik}
% ============================================================

Die Optik ($N \approx 9$) ist nicht das Ende der $\xi$-Leiter ---
aber sie ist das Ende der nutzbaren Zone für
geführte photonische Informationsübertragung.
Was darüber kommt, folgt einer klaren Abfolge:

\begin{center}
\resizebox{\textwidth}{!}{%
\begin{tabular}{p{2.5cm}p{2.5cm}p{1.5cm}p{1.5cm}p{4.5cm}}
\toprule
\textbf{Bereich} &
\textbf{Energie} &
\textbf{$N$-Stufe} &
\textbf{$\lambda$} &
\textbf{Was passiert mit dem Resonator} \\
\midrule
UV (nah) &
3--10\,eV &
$\approx 9{,}3$ &
120--40\,nm &
Photoionisation beginnt; Elektronen werden
aus äußeren Schalen geschlagen.
Resonator wird durch sein eigenes Signal
beschädigt. \\[6pt]
Weiche Röntgen &
100\,eV -- 1\,keV &
$\approx 10$ &
12--1\,nm &
$\lambda$ nähert sich dem Gitterabstand
($d \approx 0{,}3$\,nm). Bragg-Streuung
an einzelnen Gitterebenen --- kein
homogenes Medium mehr.
Wellenleitung bricht zusammen. \\[6pt]
Harte Röntgen &
1--100\,keV &
$\approx 10{,}5$ &
1--0{,}01\,nm &
$\lambda < d_\text{Gitter}$:
Photon fliegt durch Gitterlücken.
Materie teils transparent.
Kein Brechungsindex, kein Resonator. \\[6pt]
Gamma / Kernübergänge &
$>$100\,keV &
$\approx 11$ &
$<$0{,}01\,nm &
Wechselwirkung nur noch mit
Atomkernen. Völlig andere
$\xi$-Stufe --- Kernniveaus statt
Elektronenniveaus. \\[6pt]
\textbf{Paarerzeugungs-Schwelle} &
\textbf{$>$1\,MeV} &
$\approx 11{,}3$ &
--- &
\textbf{Photon erzeugt
Elektron-Positron-Paar.
Natürliche Obergrenze.} \\
\bottomrule
\end{tabular}
}
\end{center}

% ============================================================
\subsection{Die Kanonenkugel-Grenze in FFGFT:
  Skalen-Inkompatibilität zweier Torsionsmuster}
\label{sec:kanonenkugel}
% ============================================================

In der Teilchensicht ist das Argument klar:
das Photon wird zur Kanonenkugel die durch die
Gitterlücken fliegt sobald $\lambda \lesssim d_\text{Gitter}$.

In FFGFT gibt es keine Kanonenkugeln und keine Wellen
als getrennte Konzepte --- es gibt nur
\textbf{Torsionskonfigurationen des Feldes
auf verschiedenen Skalen}.

\begin{foundation}[FFGFT: Wellenleitung als Skalen-Resonanz
  zweier Torsionsmuster]

Der Kristall ist eine \textbf{periodische Torsionsstruktur}:
ein räumlich wiederholtes Muster stabiler Feldknoten
auf der $\xi$-Stufe $N \approx 7$--8 (Atomskala, $d \approx 0{,}3$\,nm).

Das Photon ist eine \textbf{propagierende Torsionswelle}:
eine kohärente Feldkonfiguration mit charakteristischer
Längenskala $\lambda$.

Wellenleitung funktioniert genau dann wenn diese
zwei Skalen kompatibel sind --- wenn die Torsionswelle
des Photons die periodische Gitterstruktur als
\emph{kollektive} homogene Torsionsumgebung wahrnimmt.
Das erfordert $\lambda \gg d_\text{Gitter}$.

\medskip
\textbf{Drei Regime in FFGFT-Sprache:}

\begin{center}
\begin{tabular}{p{3cm}p{4.5cm}p{4.5cm}}
\toprule
\textbf{Bedingung} & \textbf{Teilchenbild} & \textbf{FFGFT-Bild} \\
\midrule
$\lambda \gg d_\text{Gitter}$ &
Welle sieht homogenes Medium;
geführte Propagation &
Torsionswelle koppelt kollektiv
an Gitterknoten --- effektiver
Brechungsindex emergiert als
Mittelwert der Knotendichte \\[8pt]
$\lambda \approx d_\text{Gitter}$ &
Bragg-Streuung an einzelnen
Gitterebenen &
Torsionswelle beginnt einzelne
Gitterknoten aufzulösen ---
kollektive Kopplung bricht
in Einzelknotenkopplung auf;
nützlich für Kristallographie,
destruktiv für Wellenleitung \\[8pt]
$\lambda \ll d_\text{Gitter}$ &
Kanonenkugel fliegt durch
Gitterlücken; Materie
transparent &
Skalen-Inkompatibilität:
Torsionswelle ist zu fein um
die Gitterknoten zu ``sehen'' ---
keine Kopplung, keine Führung \\
\bottomrule
\end{tabular}
\end{center}

\medskip
\textbf{Der entscheidende Punkt:}
Teilchenbild und FFGFT-Wellenbild sind
\emph{dasselbe geometrische Argument} ---
formuliert auf verschiedenen Beschreibungsebenen.
FFGFT zeigt warum: weil es keinen fundamentalen
Unterschied zwischen Teilchen und Welle gibt.
Beide sind Torsionskonfigurationen.
Der Unterschied ist nur die \textbf{Skala der Konfiguration
relativ zur Skala der Umgebung}.

Die ``Welle-Teilchen-Dualität'' der Standard-QM
ist in FFGFT keine Dualität --- es ist eine einzige
Geometrie die je nach Beobachtungsskala
als Welle oder Teilchen erscheint.
\end{foundation}

% ============================================================
\subsection{Die Paarerzeugungs-Schwelle als natürlicher
  Endpunkt der \texorpdfstring{$\xi$}{xi}-Leiter}
\label{sec:paarerschoepfung}
% ============================================================

Die Paarerzeugungsschwelle bei $E > 2 m_e c^2 = 1{,}022$\,MeV
ist kein technisches Versagen ---
sie ist eine \textbf{geometrische Grenze der $\xi$-Hierarchie}
für Trägersysteme.

In FFGFT: das Photon ist eine propagierende
Torsionskonfiguration mit Energie $E_\text{Photon}$.
Ein Elektron und ein Positron sind zwei
\textbf{stabile elementare Torsionsknoten} ---
die fundamentalsten stabilen Konfigurationen
auf der Skala der Elementarteilchen.

Ihre Ruheenergie beträgt je $m_e c^2 = 0{,}511$\,MeV.
Zusammen: $2 m_e c^2 = 1{,}022$\,MeV.

Sobald $E_\text{Photon} \geq 2 m_e c^2$:
die Torsionsenergie des Photons übersteigt
die Schwelle um zwei neue stabile Konfigurationen
zu erzeugen.
Das Photon \emph{reorganisiert} sich in
zwei stabilere lokale Torsionsknoten ---
nicht weil es zerstört wird,
sondern weil das der energetisch mögliche
geometrische Übergang ist.

\begin{keyresult}[Paarerzeugungs-Schwelle in der $\xi$-Hierarchie]
Die Schwelle $E = 2 m_e c^2 \approx 1$\,MeV
entspricht der $\xi$-Stufe $N \approx 11{,}3$ ---
der Skala auf der Elementarteilchen als stabile
Torsionskonfigurationen existieren.

Das ist die natürliche Obergrenze für photonische
Informationsübertragung:
oberhalb dieser Schwelle ist das Photon kein
reiner Träger mehr --- es wird zum
Teilchenerzeuger.

Die nutzbare Zone der $\xi$-Hierarchie für
geführte Informationsübertragung liegt zwischen:
\begin{equation}
  N \approx 5{,}5 \;\text{(AM-Rundfunk, MHz)}
  \quad\text{bis}\quad
  N \approx 9{,}3 \;\text{(nahes UV, eV)}.
\end{equation}
Vier $\xi$-Stufen, Faktor $(1/\xi)^4 \approx 3 \times 10^{15}$
in der Frequenz --- exakt die Spanne von
1\,MHz (AM) bis $10^{15}$\,Hz (UV).
\end{keyresult}

\begin{foundation}[FFGFT: Einheitliche Sicht auf Welle,
  Teilchen und Gitter-Wechselwirkung]
In FFGFT sind drei scheinbar verschiedene Phänomene
Ausdruck derselben geometrischen Logik:

\textbf{Brechungsindex:}
Kollektive Kopplung einer Torsionswelle an
viele Gitterknoten gleichzeitig ---
emergentes Phänomen wenn $\lambda \gg d_\text{Gitter}$.

\textbf{Bragg-Streuung / Kanonenkugel-Grenze:}
Übergang von kollektiver zu Einzelknotenkopplung
wenn $\lambda \approx d_\text{Gitter}$ ---
dieselbe Skalen-Inkompatibilität
aus Teilchensicht und Wellensicht.

\textbf{Paarerzeugung:}
Torsionsenergie überschreitet die Schwelle
für neue stabile Konfigurationen ---
das Photon reorganisiert sich in
zwei Torsionsknoten (Elektron + Positron).

Alle drei sind Übergänge zwischen
Torsionskonfigurationen auf verschiedenen
$\xi$-Stufen. Die $\xi$-Hierarchie ist nicht
nur eine Längenskalen-Einteilung ---
sie ist eine Einteilung der
\textbf{stabilen Kopplungsregime}
zwischen Torsionskonfigurationen.
\end{foundation}

% ============================================================
\subsection{Konvergenz: Trägerfrequenz trifft interne
  Übergangsfrequenz auf \texorpdfstring{$\xi$}{xi}-Stufe
  \texorpdfstring{$N \approx 8$--9}{N≈8-9}}
\label{sec:konvergenz}
% ============================================================

Bisher wurden zwei Frequenzen parallel geführt,
ohne ihre Verbindung explizit zu machen:

\textbf{Trägerfrequenz:} die externe Frequenz des Lichts
das Information transportiert ---
die Frequenz die wir die $\xi$-Leiter hochklettern.

\textbf{Übergangsfrequenz:} die interne Resonanzfrequenz
eines Quantenresonators --- die Frequenz
die einem Energiesprung zwischen zwei Niveaus
entspricht: $\nu_\text{intern} = \Delta E / h$.

Diese zwei Frequenzen sind grundlegend verschieden.
Aber auf der $\xi$-Stufe $N \approx 8$--9
geschieht etwas Bemerkenswertes: sie \textbf{konvergieren}.

\medskip
\textbf{Die interne Übergangsfrequenz eines Quantenpunkts}
auf $\xi$-Stufe $N \approx 8$ (Nanometer-Skala)
liegt bei $\Delta E \approx 0{,}5$--2\,eV ---
also genau im optischen Bereich,
1550\,nm bis 700\,nm Wellenlänge.

\textbf{Die Glasfaser-Kommunikation}
arbeitet bei 1550\,nm ($E \approx 0{,}8$\,eV) ---
exakt im Fenster wo die internen Übergänge
der InGaAs-Halbleiterstrukturen liegen
die als Laser, Verstärker und Detektoren dienen.

\begin{keyresult}[Konvergenz auf $\xi$-Stufe $N \approx 8$--9]
Die Glasfaser-Kommunikation bei 1550\,nm funktioniert
so natürlich weil auf der $\xi$-Stufe $N \approx 8$--9
zwei Skalen zusammenfallen:

\begin{center}
\begin{tabular}{lll}
\toprule
& \textbf{Skala} & \textbf{Energie} \\
\midrule
Trägerfrequenz (Glasfaser) & $\lambda = 1550$\,nm & 0{,}8\,eV \\
Interne Übergänge (InGaAs-QD) & $R \approx 3$--5\,nm & 0{,}7--1{,}2\,eV \\
Exziton-Bohr-Radius (CdSe) & $a_X \approx 5{,}5$\,nm & $\sim$1\,eV \\
\bottomrule
\end{tabular}
\end{center}

Das ist keine Materialkunde-Konvention ---
es ist geometrisch erzwungen durch die $\xi$-Hierarchie.
Dieselbe Skalenstufe die den Brechungsindex
des Wellenleiters bestimmt bestimmt auch
die internen Energieniveaus der aktiven Komponenten.
Träger und Resonator sprechen dieselbe Sprache
weil sie auf derselben $\xi$-Stufe leben.

Und deshalb bricht die Wellenleitung oberhalb
des optischen Fensters nicht nur aus Gittergeometrie
zusammen --- sondern auch weil die Trägerfrequenz
die $\xi$-Stufe verlässt auf der die aktiven
Materialkomponenten operieren.
\end{keyresult}

% ============================================================
\section{\texorpdfstring{$\xi$}{xi}-Stufentrennung als Stabilitätsprinzip
  der Naturhierarchie}
\label{sec:stufentrennung}
% ============================================================

Die $\xi$-Hierarchie ist mehr als eine Einteilung
von Längenskalen. Sie ist eine
\textbf{Stabilitätsstruktur}.

Jede Stufe $N$ ist ein stabiles Kopplungsregime ---
Torsionskonfigurationen auf dieser Skala koppeln
stark aneinander (innerhalb der Stufe)
und schwach an alle anderen Stufen
(zwischen den Stufen).

% ============================================================
\subsection{Die Kopplungsformel zwischen Stufen}
% ============================================================

Die Stärke der Kopplung zwischen zwei Torsions\-konfigurationen
auf den Stufen $N_1$ und $N_2$ skaliert mit:
\begin{equation}
  \boxed{C(N_1, N_2) \;\propto\; \xi^{|N_1 - N_2|}.}
\end{equation}
Da $\xi \approx 1{,}333 \times 10^{-4} \ll 1$,
fällt die Kopplung exponentiell mit dem Stufenabstand:

\begin{center}
\resizebox{\textwidth}{!}{%
\begin{tabular}{rrrl}
\toprule
$\Delta N = |N_1 - N_2|$ &
$\xi^{\Delta N}$ &
Größenordnung &
Konsequenz \\
\midrule
0 & 1 & $10^0$ &
Volle Kopplung --- gleiche Stufe,
starke Wechselwirkung \\[4pt]
1 & $\xi$ & $10^{-4}$ &
Schwache Kopplung ---
z.\,B. Photon koppelt an
Quantenpunkt-Übergang \\[4pt]
2 & $\xi^2$ & $10^{-8}$ &
Sehr schwach ---
z.\,B. optisches Photon
an Kernniveau \\[4pt]
3 & $\xi^3$ & $10^{-12}$ &
Praktisch null ---
z.\,B. Phonon (Atom-Stufe)
an Molekülstruktur \\[4pt]
4 & $\xi^4$ & $10^{-16}$ &
Nicht messbar ---
z.\,B. atomare Schwingung
an zelluläre Struktur \\[4pt]
6 & $\xi^6$ & $10^{-24}$ &
Absolut vernachlässigbar ---
z.\,B. atomare Schwingung
an kosmologische Struktur \\
\bottomrule
\end{tabular}
}
\end{center}

% ============================================================
\subsection{Fraktale Selbstähnlichkeit \emph{wegen} der Dämpfung}
% ============================================================

Das ist der entscheidende Punkt:

\textbf{Die fraktale Hierarchie der Natur ist nicht trotz
der $\xi$-Dämpfung stabil --- sie ist stabil
\emph{wegen} ihr.}

Wäre $\xi$ nicht klein, würden alle Stufen
stark aneinander koppeln.
Jede Änderung auf einer Stufe würde sofort
alle anderen Stufen umstrukturieren.
Es gäbe keine stabilen Atome, keine stabilen Moleküle,
keine stabilen Zellen --- nur ein einziges
chaotisches gekoppeltes System aller Skalen gleichzeitig.

Weil $\xi \ll 1$ ist, können die Stufen
\textbf{quasi-unabhängig} existieren:
jede Stufe hat ihre eigene stabile interne Dynamik,
die von den anderen Stufen kaum gestört wird.

\begin{foundation}[FFGFT: $\xi$-Dämpfung als Ursprung der
  Naturhierarchie]
Die Existenz getrennter, stabiler Strukturebenen
in der Natur --- Quarks, Nukleonen, Atome,
Moleküle, Zellen, Organismen, Planeten, Galaxien ---
ist keine zufällige Eigenschaft der Materie.

Sie folgt direkt aus der $\xi$-Stufentrennung:
Jede Ebene ist für die übernächste praktisch unsichtbar.
Die Kopplung fällt exponentiell mit dem Stufenabstand.
Das ist der Mechanismus der die fraktale
Selbstähnlichkeit der Natur \emph{stabil} macht.

\textbf{Die seltenen Ausnahmen} --- Stellen wo eine Stufe
die andere tatsächlich beeinflusst ---
sind geometrisch vorhersagbar:
es sind immer Fälle wo $\Delta N$ klein ist,
also die Energieskalen nahe beieinander liegen:

\begin{center}
\begin{tabular}{p{4cm}cp{6cm}}
\toprule
\textbf{Phänomen} &
\textbf{$\Delta N$} &
\textbf{Warum Kopplung möglich} \\
\midrule
Quantenpunkt absorbiert
optisches Photon &
$\approx 1$ &
$\xi^1 \approx 10^{-4}$: schwach,
aber resonant verstärkt ---
Energieniveaus stimmen überein \\[6pt]
Photosynthese: Photon
regt FMO-Komplex an &
$\approx 1$ &
Proteinhülle justiert Niveaus
auf $\Delta N \approx 1$ ---
biologische Feinabstimmung
der $\xi$-Resonanz \\[6pt]
Phonon stört
Spin-Kohärenz (Dekohärenz) &
$\approx 0{,}5$ &
Phonon-Energie ($\sim$meV)
und Spin-Aufspaltung ($\sim$meV)
auf gleicher Stufe ---
deshalb Dekohärenz unvermeidlich \\[6pt]
Laserpuls schaltet
Quantenpunkt-Zustand &
$\approx 1$ &
Resonante Kopplung über
$\Delta N \approx 1$ ---
gezielt angesteuert \\[6pt]
Optisches Photon passiert
Gitter ohne Streuung &
$\approx 2$ &
$\xi^2 \approx 10^{-8}$:
zu schwach für Kopplung ---
Materie transparent \\
\bottomrule
\end{tabular}
\end{center}
\end{foundation}

% ============================================================
\subsection{Umstrukturierung zwischen Stufen:
  wann passiert sie und wann nicht?}
% ============================================================

Die Frage ``kann eine Stufe die andere umstrukturieren?''
hat jetzt eine präzise Antwort:

\textbf{Ja, wenn:}
\begin{itemize}
  \item $\Delta N$ klein ist ($\leq 1$) --- die Energieskalen
    sind nah genug dass die Kopplung $\xi^{\Delta N}$
    durch resonante Verstärkung kompensiert werden kann.
  \item Oder die Energie auf Stufe $N_1$ die
    \emph{Schwelle} für eine neue stabile Konfiguration
    auf Stufe $N_2$ übersteigt ---
    wie bei der Paarerzeugung ($\gamma \to e^+ + e^-$).
\end{itemize}

\textbf{Nein, wenn:}
\begin{itemize}
  \item $\Delta N \geq 2$ --- die Kopplung
    $\xi^{\Delta N} \leq 10^{-8}$ ist zu schwach
    um stabile Umstrukturierung zu verursachen.
  \item Die Energie reicht nicht aus die Schwelle
    der Zielkonfiguration zu überwinden.
\end{itemize}

\begin{keypoint}[Die Naturhierarchie als geometrisches
  Stabilitätssystem]
Atome werden nicht durch Phononen umstrukturiert
($\Delta N \approx 1$, Kopplung $\sim \xi$: schwach,
aber nicht null --- deshalb Dekohärenz existiert,
aber Atome trotzdem stabil bleiben).

Zellen werden nicht durch atomare Schwingungen
umstrukturiert ($\Delta N \approx 4$,
Kopplung $\sim \xi^4 \approx 10^{-16}$:
absolut vernachlässigbar).

Galaxien werden nicht durch Phononen umstrukturiert
($\Delta N \approx 8$,
Kopplung $\sim \xi^8 \approx 10^{-32}$:
nicht existent).

Die Hierarchie der Natur --- von Quarks bis Galaxien ---
ist in FFGFT kein Zufall der Materieverteilung.
Sie ist die direkte Folge der $\xi$-Stufentrennung:
jede Ebene ist geometrisch gegen Störungen
von weit entfernten Stufen geschützt.
Selbstähnlichkeit und Stabilität sind
zwei Seiten derselben $\xi$-Geometrie.
\end{keypoint}

% ============================================================

\section{Qudit-Register --- die Verbindung zur \texorpdfstring{$\xi$}{xi}-Hierarchie}
\label{sec:qudit-register}
% ============================================================

Aus Kapitel~174: ein Quantenpunkt der Dimension $d$
kodiert $\log_2 d$ Bits.
Was passiert wenn man $n$ solche Qudits verschränkt?

Der Gesamtzustandsraum hat Dimension $d^n$.
Für $d = 4$ (Spin-Orbital-Ququart, zwei unterste Orbitale
$\times$ Spin-up/down) und $n$ Qudits:
\begin{equation}
  d^n \;=\; 4^n \;=\; (2^2)^n \;=\; 2^{2n}.
\end{equation}
Das ist äquivalent zu $2n$ Qubits in einem Register
der Größe $n$ --- doppelte Informationsdichte pro Objekt.

\begin{keyresult}[Qudit-Dichte auf der $\xi$-Stufe $N \approx 8$]
Ein Quantenpunkt mit $R = 3$\,nm bei 4\,K
hat mindestens 4 thermisch stabile Zustände
(Spin-Orbital-Ququart, $d = 4$).
Ein Register von $n = 10$ solchen Qudits
hat Dimension $4^{10} = 1{,}048{,}576$ ---
dasselbe wie 20 klassische Qubits, aber
nur 10 physische Objekte auf der $\xi$-Stufe $N \approx 8$.

Die $\xi$-Hierarchie erzwingt diese Stufe:
Alle diese Quantenpunkte sitzen geometrisch
bei $N \approx 8$. Das ist kein Zufall der
Materialkunde --- es ist die geometrische Grundlage
(Kapitel~173), die vorschreibt welche Skalen
stabile Quantenresonatoren ausbilden können.
\end{keyresult}

% ============================================================
\section{Der T0-Topologie-Compiler --- Geometrie der Verschränkung}
\label{sec:compiler}
% ============================================================

Ein praktisches Problem bei verschränkten Registern:
Verschränkungs-Gatter zwischen nicht-benachbarten Qubits
sind teuer --- sie erfordern viele SWAP-Operationen
(die Fehler akkumulieren und Zeit kosten).

Die T0-Theorie (Dok.~034) bietet einen geometrischen
Lösungsweg: den \textbf{T0-Topologie-Compiler}.

Standardkompiler minimieren die Zahl der SWAPs.
Der T0-Kompiler minimiert stattdessen die kumulative
\textbf{fraktale Weglänge} aller Verschränkungsoperationen.
Da der fraktale Dämpfungsterm $\exp(-\xi \cdot \ell / D_f)$
abstandsabhängig ist (größere physische Distanz $\ell$
bedeutet stärkere $\xi$-Dämpfung der Kohärenz),
müssen kritisch interagierende Qubits
\emph{physisch näher angeordnet} werden.

Das führt zu einer neuen Chip-Architektur:
nicht nach logischer Schaltkreisoptimierung ausgelegt,
sondern nach \textbf{geometrischer Kohärenzmaximierung}.

\begin{significance}[Verbindung zur Adaptiven Hybridarchitektur]
Die adaptive Hybridarchitektur (Kapitel~173-Schluss)
wählt Resonatoren je nach Aufgabe:
Quantenpunkte für lokale synaptische Operationen,
Ising-Maschinen für globale Optimierung.

Der T0-Topologie-Compiler ist das Bindeglied:
Er entscheidet welche Qubits/Qudits physisch
nahe beieinander sitzen müssen (starke Verschränkung)
und welche weiter entfernt bleiben dürfen
(schwache, tolerant gedämpfte Verschränkung).

Die lernende KI auf dieser Hardware lernt nicht nur
\emph{was} berechnet werden soll ---
sie lernt auch \emph{welche geometrische Anordnung}
die Berechnung am kohärentesten ausführt.
Nicht durch programmierte Topologie-Vorgaben,
sondern durch geometrisch informiertes Lernen
im T0-Kompiler.
\end{significance}

% ============================================================
\section{Vergleich: klassisches Bit, Qubit, Qudit, Register}
\label{sec:vergleich}
% ============================================================

\begin{center}
\resizebox{\textwidth}{!}{%
\begin{tabular}{p{2.5cm}p{2.5cm}p{2.5cm}p{2.5cm}p{2.5cm}}
\toprule
& \textbf{Klass. Bit} & \textbf{Qubit} &
\textbf{Qudit ($d$)} & \textbf{$n$ Qubits} \\
\midrule
\textbf{Zustandsraum} &
2 Punkte &
Bloch-Kugel ($\infty$) &
$\mathbb{CP}^{d-1}$ ($\infty$) &
$\mathbb{CP}^{2^n-1}$ ($\infty$) \\[4pt]
\textbf{Auslese-Ergebnis} &
1 Bit (deterministisch) &
1 Bit (deterministisch,\newline abhängig vom Phasenzustand) &
$\log_2 d$ Bits (deterministisch) &
$n$ \\[4pt]
\textbf{Parallelverarbeitung} &
keine &
Interferenz über 2 Pfade &
Interferenz über $d$ Pfade &
Interferenz über $2^n$ Pfade \\[4pt]
\textbf{Ressource für Vorteil} &
--- &
Phase $\phi$ &
$d-1$ Phasen &
Verschränkung + Phasen \\[4pt]
\textbf{FFGFT-Bild} &
stabile/instabile Torsion &
Torsion auf Bloch-Kugel &
$d$ stabile Torsionsniveaus &
Torsionsnetz, topol. verbunden \\
\bottomrule
\end{tabular}
}
\end{center}

% ============================================================
\section{Offene Verbindungen --- nächste Schritte}
\label{sec:offene-175}
% ============================================================

\begin{significance}[Offene Fragen für weitere Kapitel]
\textbf{1. Fraktale Gatter-Korrekturen messbar?}\\
Der T0-Kompiler sagt: alle $\pi$-Pulse weichen
um $\Delta\alpha \approx \Kfrak - 1 \approx 0{,}013\,\%$
vom idealen $\pi$ ab.
Heutige supraleitende Qubit-Systeme erreichen
Gatter-Treue von $>99{,}9\,\%$.
Das ist gerade an der Grenze um diese
systematische Abweichung zu sehen.
Gibt es in aktuellen Kalibrierungsdaten
einen systematischen Bias in der $\pi$-Puls-Länge?\\[6pt]

\textbf{2. Bell-Zustände und $\xi$-Skalenstufe:}\\
Die $\xi$-Dämpfung in Bell-Korrelationen
($\Delta\text{CHSH} \approx 10^{-4}$) hängt von $N \approx 8$ ab.
Für Bell-Tests mit \emph{makroskopischen} Objekten
(Atome, Moleküle) würde die Skalenstufe $N$ steigen ---
und damit die Dämpfung zunehmen.
Gibt es eine $N$-abhängige Unterdrückung
von Bell-Verletzungen für verschiedene Systemgrößen?\\[6pt]

\textbf{3. Qudit-Verschränkung und Torsustopologie:}\\
Zwei verschränkte Qudits ($d = 4$) entsprechen
in FFGFT zwei topologisch verbundenen
Spin-Orbital-Torsionskonfigurationen.
Was ist die geometrische Beschreibung
dieser Verbindung im toroidalen Feld?
Gibt es eine Erhaltungsgröße (topologische Invariante)
die die Verschränkungstiefe quantifiziert?\\[6pt]

\textbf{4. Quantenfehlerkorrektur und $\xi$-Granularität:}\\
Der Surface-Code braucht $\sim 1000$ physische Qubits
pro logischem Qubit.
In einem Qudit-basierten Register mit $d = 4$
braucht man nur halb so viele physische Objekte
für dieselbe logische Information.
Wie verändert sich die Fehlerkorrektur-Schwelle
im T0-Rahmen wenn die natürliche Rauschgrenze
$\xi \approx 1{,}333 \times 10^{-4}$ als
untere Fehlerschranke gilt?
\end{significance}

% ============================================================
\section{Zusammenfassung}
\label{sec:zusammenfassung-175}
% ============================================================

\begin{conclusionBox}[Fazit: Mehr als zwei Zustände]
Ein einzelnes Qubit hat unendlich viele Zwischenzustände
im Torsionsphasenraum --- die Zustände selbst sind aber
\textbf{diskret}: das Auslesen liefert immer einen der
zwei stabilen Pole, deterministisch bestimmt durch den
genauen Phasenzustand zum Auslesezeitpunkt.
Die scheinbare ``Wahrscheinlichkeit'' der Standard-QM
ist epistemische Unkenntnis dieses Zustands ---
keine fundamentale Zufälligkeit des Systems.

\textbf{Das Auslesen wechselwirkt mit dem System ---
aber nicht immer störend.}
Die Faraday/Kerr-Rotation (Kapitel~174) ist eine
QND-Messung: sie liest den Zustand aus ohne ihn zu verändern,
weil die Kopplung mit dem Spin-Operator kommutiert.
Dass derselbe Zustand wiederholbar gemessen werden kann,
ist der empirische Beweis dass er \emph{vor} der Messung
existiert hat --- definitiv und real.
Deterministisch bleibt das Ergebnis in jedem Fall.

Der eigentliche Quantengewinn kommt auf drei Wegen:

\textbf{1. Superposition:}
Kontinuierliche intermediäre Torsion mit Phase $\phi$
ermöglicht Interferenz --- die geometrische Grundlage
aller nützlichen Quantenalgorithmen.
Interferenz richtet das Register so aus, dass
die gesuchte Antwort in einer stabilen Torsionskonfiguration
endet; das Auslesen folgt dem deterministisch.

\textbf{2. Register-Verschränkung:}
$n$ Qubits spannen einen $2^n$-dimensionalen Torsionsraum auf.
Verschränkung bedeutet: die Torsionsknoten der Qubits
sind topologisch verbunden --- das Auslesen eines
Knotens bestimmt die anderen zwingend mit.
Bell-Korrelationen werden durch $\xi$-Dämpfung
lokal erklärt ohne Nicht-Lokalität zu beseitigen.

\textbf{3. Qudit-Kodierung:}
Ein Resonator mit $d > 2$ stabilen Torsionsniveaus
liefert beim Auslesen $\log_2 d > 1$ Bits pro Objekt.
Spin-Orbital-Ququarts ($d = 4$) auf der
$\xi$-Stufe $N \approx 8$ verdoppeln die
Informationsdichte eines Registers.

In FFGFT sind alle drei Strategien Ausdruck desselben
Grundprinzips: stabile Torsionskonfigurationen
in geometrisch präzisen Resonatoren --- diskret von innen,
analog steuerbar von außen, deterministisch beim Auslesen.
\end{conclusionBox}

% ============================================================
% ============================================================

\begin{thebibliography}{99}

% --- Quantenpunkt-Grundlagen ---
\bibitem{Kouwenhoven1997}
Kouwenhoven, L.\,P. et al.
\textit{Excitation spectra of circular, few-electron quantum dots.}
Science \textbf{278}, 1788--1792 (1997).

% --- Loss-DiVincenzo ---
\bibitem{Loss1998}
Loss, D. \& DiVincenzo, D.\,P.
\textit{Quantum computation with quantum dots.}
Phys.\ Rev.\ A \textbf{57}, 120--126 (1998).

% --- Spin-Qubit Kohärenz in Si ---
\bibitem{Veldhorst2014}
Veldhorst, M. et al.
\textit{An addressable quantum dot qubit with fault-tolerant control-fidelity.}
Nature Nanotechnol.\ \textbf{9}, 981--985 (2014).

\bibitem{Muhonen2014}
Muhonen, J.\,T. et al.
\textit{Storing quantum information for 30 seconds in a nanoelectronic device.}
Nature Nanotechnol.\ \textbf{9}, 986--991 (2014).

% --- Optische QND-Auslese ---
\bibitem{Atature2007}
Atatüre, M. et al.
\textit{Observation of Faraday rotation from a single confined spin.}
Nature Phys.\ \textbf{3}, 101--105 (2007).

\bibitem{Grangier1998}
Grangier, P., Levenson, J.\,A. \& Poizat, J.\,P.
\textit{Quantum non-demolition measurements in optics.}
Nature \textbf{396}, 537--542 (1998).

% --- Photonische Wellenleiter und LiNbO3 ---
\bibitem{Wang2018}
Wang, C. et al.
\textit{Integrated lithium niobate electro-optic modulators operating at
CMOS-compatible voltages.}
Nature \textbf{562}, 101--104 (2018).

\bibitem{Zhu2021}
Zhu, D. et al.
\textit{Integrated photonics on thin-film lithium niobate.}
Adv.\ Opt.\ Photon.\ \textbf{13}, 242--352 (2021).

% --- MZI als QFT / Unitäre Transformation ---
\bibitem{Reck1994}
Reck, M. et al.
\textit{Experimental realization of any discrete unitary operator.}
Phys.\ Rev.\ Lett.\ \textbf{73}, 58--61 (1994).

% --- Shannon-Kapazität ---
\bibitem{Shannon1948}
Shannon, C.\,E.
\textit{A mathematical theory of communication.}
Bell Syst.\ Tech.\ J.\ \textbf{27}, 379--423 (1948).

% --- Bragg-Streuung und Röntgenbeugung ---
\bibitem{Bragg1913}
Bragg, W.\,H. \& Bragg, W.\,L.
\textit{The reflection of X-rays by crystals.}
Proc.\ R.\ Soc.\ A \textbf{88}, 428--438 (1913).

% --- Paarerzeugung ---
\bibitem{Dirac1928}
Dirac, P.\,A.\,M.
\textit{The quantum theory of the electron.}
Proc.\ R.\ Soc.\ A \textbf{117}, 610--624 (1928).

% --- Dekohärenz durch Kernspins ---
\bibitem{Khaetskii2002}
Khaetskii, A.\,V., Loss, D. \& Glazman, L.
\textit{Electron spin decoherence in quantum dots due to interaction with
nuclei.}
Phys.\ Rev.\ Lett.\ \textbf{88}, 186802 (2002).

% --- KLM-Protokoll (für Kapitel 175 Abschnitt Shor/photonisch) ---
\bibitem{Knill2001}
Knill, E., Laflamme, R. \& Milburn, G.\,J.
\textit{A scheme for efficient quantum computation with linear optics.}
Nature \textbf{409}, 46--52 (2001).

\bibitem{Kok2007}
Kok, P. et al.
\textit{Linear optical quantum computing with photonic qubits.}
Rev.\ Mod.\ Phys.\ \textbf{79}, 135--174 (2007).

\end{thebibliography}
